\begin{thebibliography}{99}
\bibitem{Akin}E. Akin - The general topology of dynamical systems, Graduate Studies in Math., vol. 1, American Mathematical Society, Providence, RI, 1993 | MR | Zbl
\bibitem{AY}V. M. Alekseev \& M. V. Yakobson - “Symbolic dynamics and hyperbolic dynamic systems”, Phys. Rep. 75 (1981) no. 5, p. 287-325 | MR | DOI
\bibitem{Anosov}D. V. Anosov - “Roughness of geodesic flows on compact Riemannian manifolds of negative curvature”, Dokl. Akad. Nauk SSSR 145 (1962), p. 707-709 | MR | Zbl
\bibitem{AH}N. Aoki \& K. Hiraide - Topological theory of dynamical systems. Recent advances, North-Holland Math. Library, vol. 52, North-Holland Publishing Co., Amsterdam, 1994 | Zbl
\bibitem{Arnold}L. Arnold - Random dynamical systems, Springer Monographs in Math., Springer-Verlag, Berlin, 1998 | Zbl | DOI
\bibitem{AB}A. Avila \& J. Bochi - “A uniform dichotomy for generic $\mathrm{SL}(2,\mathbb{R})$ cocycles over a minimal base”, Bull. Soc. math. France 135 (2007) no. 3, p. 407-417 | MR | Zbl | DOI
\bibitem{AV_Portugalia}A. Avila \& M. Viana - “Simplicity of Lyapunov spectra: a sufficient criterion”, Portugal. Math. 64 (2007) no. 3, p. 311-376 | MR | Zbl | DOI
\bibitem{AV_Inventiones}A. Avila \& M. Viana - “Extremal Lyapunov exponents: an invariance principle and applications”, Invent. Math. 181 (2010) no. 1, p. 115-189 | MR | Zbl | DOI
\bibitem{Ball}K. Ball - “An elementary introduction to modern convex geometry”, in Flavors of geometry, Math. Sci. Res. Inst. Publ., vol. 31, Cambridge Univ. Press, Cambridge, 1997, p. 1-58 | MR | Zbl
\bibitem{Barabanov}N. E. Barabanov - “On the Lyapunov exponent of discrete inclusions. I”, Avtomat. i Telemekh. (1988) no. 2, p. 40-46 | MR | Zbl
\bibitem{BWang}M. A. Berger \& Y. Wang - “Bounded semigroups of matrices”, Linear Algebra Appl. 166 (1992), p. 21-27 | MR | Zbl | DOI
\bibitem{BT}V. D. Blondel \& J. N. Tsitsiklis - “The boundedness of all products of a pair of matrices is undecidable”, Systems Control Lett. 41 (2000) no. 2, p. 135-140 | MR | Zbl | DOI
\bibitem{B_Oseledets}J. Bochi - “The multiplicative ergodic theorem of Oseledets”, 2008, Note available at http://www.mat.uc.cl/~jairo.bochi/docs/oseledets.pdf
\bibitem{Bochi_ICM}J. Bochi - “Ergodic optimization of Birkhoff averages and Lyapunov exponents”, in Proc. Int. Cong. of Math. (Rio de Janeiro, 2018), Vol. 2, World Scientific, 2019, p. 1821-1842
\bibitem{BG}J. Bochi \& N. Gourmelon - “Some characterizations of domination”, Math. Z. 263 (2009) no. 1, p. 221-231 | MR | Zbl | DOI
\bibitem{BM}J. Bochi \& I. D. Morris - “Continuity properties of the lower spectral radius”, Proc. London Math. Soc. (3) 110 (2015) no. 2, p. 477-509 | MR | Zbl | DOI
\bibitem{BPS}J. Bochi, R. Potrie \& A. Sambarino - “Anosov representations and dominated splittings”, J. Eur. Math. Soc. (JEMS) 21 (2019) no. 11, p. 3343-3414 | MR | DOI
\bibitem{BR}J. Bochi \& M. Rams - “The entropy of Lyapunov-optimizing measures of some matrix cocycles”, J. Modern Dyn. 10 (2016), p. 255-286 | MR | Zbl | DOI
\bibitem{BV}C. Bonatti \& M. Viana - “Lyapunov exponents with multiplicity 1 for deterministic products of matrices”, Ergodic Theory Dynam. Systems 24 (2004) no. 5, p. 1295-1330 | MR | Zbl | DOI
\bibitem{BDV_book}C. Bonatti, L. J. Díaz \& M. Viana - Dynamics beyond uniform hyperbolicity. A global geometric and probabilistic perspective. Mathematical Physics, III, Encyclopaedia of Math. Sciences, vol. 102, Springer-Verlag, Berlin, 2005 | Zbl
\bibitem{BGV}C. Bonatti, X. Gómez-Mont \& M. Viana - “Généricité d’exposants de Lyapunov non-nuls pour des produits déterministes de matrices”, Ann. Inst. H. Poincaré Anal. Non Linéaire 20 (2003) no. 4, p. 579-624 | Zbl | DOI
\bibitem{Bousch_Walters}T. Bousch - “La condition de Walters”, Ann. Sci. École Norm. Sup. (4) 34 (2001) no. 2, p. 287-311 | MR | Zbl | DOI
\bibitem{Bousch_amphi}T. Bousch - “Le lemme de Mañé-Conze-Guivarc’h pour les systèmes amphidynamiques rectifiables”, Ann. Fac. Sci. Toulouse Math. (6) 20 (2011) no. 1, p. 1-14 | MR | Zbl | DOI
\bibitem{Bousch_Jenkinson}T. Bousch \& O. Jenkinson - “Cohomology classes of dynamically non-negative $C^k$ functions”, Invent. Math. 148 (2002) no. 1, p. 207-217 | MR | Zbl | DOI
\bibitem{Bousch_Mairesse}T. Bousch \& J. Mairesse - “Asymptotic height optimization for topical IFS, Tetris heaps, and the finiteness conjecture”, J. Amer. Math. Soc. 15 (2002) no. 1, p. 77-111 | MR | Zbl | DOI
\bibitem{Bowen}R. Bowen - “Periodic points and measures for Axiom $A$ diffeomorphisms”, Trans. Amer. Math. Soc. 154 (1971), p. 377-397 | MR | Zbl | DOI
\bibitem{BQuas}X. Bressaud \& A. Quas - “Rate of approximation of minimizing measures”, Nonlinearity 20 (2007) no. 4, p. 845-853 | MR | Zbl | DOI
\bibitem{BreuF}E. Breuillard \& K. Fujiwara - “On the joint spectral radius for isometries of non-positively curved spaces and uniform growth”, 2018 | arXiv
\bibitem{BrinP}M. I. Brin \& J. B. Pesin - “Partially hyperbolic dynamical systems”, Izv. Akad. Nauk SSSR Ser. Mat. 38 (1974), p. 170-212 | MR | Zbl
\bibitem{CGP}A. Cicone, N. Guglielmi \& V. Y. Protasov - “Linear switched dynamical systems on graphs”, Nonlinear Anal. Hybrid Syst. 29 (2018), p. 165-186 | MR | Zbl | DOI
\bibitem{CK_book}F. Colonius \& W. Kliemann - The dynamics of control, Systems \& Control: Foundations \& Applications, Birkhäuser Boston, Inc., Boston, MA, 2000 | Zbl | DOI
\bibitem{Contreras}G. Contreras - “Ground states are generically a periodic orbit”, Invent. Math. 205 (2016) no. 2, p. 383-412 | MR | Zbl | DOI
\bibitem{CLT}G. Contreras, A. O. Lopes \& P. Thieullen - “Lyapunov minimizing measures for expanding maps of the circle”, Ergodic Theory Dynam. Systems 21 (2001) no. 5, p. 1379-1409 | MR | Zbl | DOI
\bibitem{CG}J. Conze \& Y. Guivarc’h - “Croissance des sommes ergodiques et principe variationnel”, circa 1993, Unpublished manuscript
\bibitem{CNP}D. Coronel, A. Navas \& M. Ponce - “On bounded cocycles of isometries over minimal dynamics”, J. Modern Dyn. 7 (2013) no. 1, p. 45-74 | MR | Zbl | DOI
\bibitem{CroPo}S. Crovisier \& R. Potrie - “Introduction to partially hyperbolic dynamics”, 2015, Notes ICTP
\bibitem{Garibaldi_book}E. Garibaldi - Ergodic optimization in the expanding case. Concepts, tools and applications, SpringerBriefs in Math., Springer, Cham, 2017 | Zbl | DOI
\bibitem{GG}E. Garibaldi \& J. T. A. Gomes - “Aubry set for asymptotically sub-additive potentials”, Stochastic Dyn. 16 (2016) no. 2, article ID 1660009, 13 pages | MR | Zbl | DOI
\bibitem{Gourmelon}N. Gourmelon - “Adapted metrics for dominated splittings”, Ergodic Theory Dynam. Systems 27 (2007) no. 6, p. 1839-1849 | MR | Zbl | DOI
\bibitem{Herman}M.-R. Herman - “Une méthode pour minorer les exposants de Lyapounov et quelques exemples montrant le caractère local d’un théorème d’Arnol’d et de Moser sur le tore de dimension $2$”, Comment. Math. Helv. 58 (1983) no. 3, p. 453-502 | MR | Zbl | DOI
\bibitem{HPS}M. W. Hirsch, C. C. Pugh \& M. Shub - Invariant manifolds, Lect. Notes in Math., vol. 583, Springer-Verlag, Berlin-New York, 1977 | Zbl
\bibitem{Jenkinson_survey}O. Jenkinson - “Ergodic optimization”, Discrete Contin. Dynam. Systems 15 (2006) no. 1, p. 197-224 | MR | Zbl | DOI
\bibitem{Jenkinson_survey_new}O. Jenkinson - “Ergodic optimization in dynamical systems”, Ergodic Theory Dynam. Systems 39 (2019) no. 10, p. 2593-2618 | MR | Zbl | DOI
\bibitem{Jungers}R. Jungers - The joint spectral radius. Theory and applications, Lect. Notes in Control and Information Sci., vol. 385, Springer-Verlag, Berlin, 2009 | MR | DOI
\bibitem{Kalinin}B. Kalinin - “Livšic theorem for matrix cocycles”, Ann. of Math. (2) 173 (2011) no. 2, p. 1025-1042 | MR | Zbl | DOI
\bibitem{KalSad}B. Kalinin \& V. Sadovskaya - “Cocycles with one exponent over partially hyperbolic systems”, Geom. Dedicata 167 (2013), p. 167-188 | MR | Zbl | DOI
\bibitem{KH}A. Katok \& B. Hasselblatt - Introduction to the modern theory of dynamical systems, Encyclopedia of Math. and its Appl., vol. 54, Cambridge University Press, Cambridge, 1995 | MR | Zbl | DOI
\bibitem{Katz}Y. Katznelson - An introduction to harmonic analysis, Dover Publications, Inc., New York, 1976 | Zbl
\bibitem{Kozyakin}V. Kozyakin - “An explicit Lipschitz constant for the joint spectral radius”, Linear Algebra Appl. 433 (2010) no. 1, p. 12-18 | MR | Zbl | DOI
\bibitem{Krengel}U. Krengel - Ergodic theorems, De Gruyter Studies in Math., vol. 6, Walter de Gruyter \& Co., Berlin, 1985 | MR | Zbl | DOI
\bibitem{LopesT}A. O. Lopes \& P. Thieullen - “Sub-actions for Anosov diffeomorphisms”, in Geometric methods in dynamics. II, Astérisque, vol. 287, Société Mathématique de France, Paris, 2003, p. 135-146 | Zbl
\bibitem{Mat}J. N. Mather - “Action minimizing invariant measures for positive definite Lagrangian systems”, Math. Z. 207 (1991) no. 2, p. 169-207 | MR | Zbl | DOI
\bibitem{Morris_sufficient}I. D. Morris - “A sufficient condition for the subordination principle in ergodic optimization”, Bull. London Math. Soc. 39 (2007) no. 2, p. 214-220 | MR | Zbl | DOI
\bibitem{Morris_rapidly}I. D. Morris - “A rapidly-converging lower bound for the joint spectral radius via multiplicative ergodic theory”, Adv. Math. 225 (2010) no. 6, p. 3425-3445 | MR | Zbl | DOI
\bibitem{Morris_Mather}I. D. Morris - “Mather sets for sequences of matrices and applications to the study of joint spectral radii”, Proc. London Math. Soc. (3) 107 (2013) no. 1, p. 121-150 | MR | Zbl | DOI
\bibitem{Oregon}E. Oregón-Reyes - “A new inequality about matrix products and a Berger-Wang formula”, 2017 | arXiv
\bibitem{Pesin_book}Y. B. Pesin - Lectures on partial hyperbolicity and stable ergodicity, Zurich Lectures in Advanced Math., European Mathematical Society, Zürich, 2004 | MR | Zbl | DOI
\bibitem{PEDJ}M. Philippe, R. Essick, G. E. Dullerud \& R. M. Jungers - “Stability of discrete-time switching systems with constrained switching sequences”, Automatica J. IFAC 72 (2016), p. 242-250 | MR | Zbl | DOI
\bibitem{PintoRand}A. A. Pinto \& D. A. Rand - “Smoothness of holonomies for codimension 1 hyperbolic dynamics”, Bull. London Math. Soc. 34 (2002) no. 3, p. 341-352 | MR | Zbl | DOI
\bibitem{PSW1}C. Pugh, M. Shub \& A. Wilkinson - “Hölder foliations”, Duke Math. J. 86 (1997) no. 3, p. 517-546 | Zbl | DOI
\bibitem{PSW2}C. Pugh, M. Shub \& A. Wilkinson - “Hölder foliations, revisited”, J. Modern Dyn. 6 (2012) no. 1, p. 79-120 | Zbl | DOI
\bibitem{Pugh}C. C. Pugh - “On arbitrary sequences of isomorphisms in $R^{m}\rightarrow R^{m}$”, Trans. Amer. Math. Soc. 184 (1973), p. 387-400 | MR | DOI
\bibitem{QZL}L. Qiu, Y. Zhang \& C.-K. Li - “Unitarily invariant metrics on the Grassmann space”, SIAM J. Matrix Anal. Appl. 27 (2005) no. 2, p. 507-531 | MR | Zbl | DOI
\bibitem{RS}G.-C. Rota \& G. Strang - “A note on the joint spectral radius”, Nederl. Akad. Wetensch. Indag. Math. 22 (1960), p. 379-381 | MR | Zbl | DOI
\bibitem{Ruelle}D. Ruelle - Thermodynamic formalism. The mathematical structures of classical equilibrium statistical mechanics, Encyclopedia of Math. and its Appl., vol. 5, Addison-Wesley Publishing Co., Reading, Mass., 1978 | Zbl
\bibitem{Sakai}K. Sakai - “Shadowing properties of $\mathcal{L}$-hyperbolic homeomorphisms”, Topology Appl. 112 (2001) no. 3, p. 229-243 | MR | Zbl | DOI
\bibitem{Sav}S. V. Savchenko - “Homological inequalities for finite topological Markov chains”, Funkcional. Anal. i Priložen. 33 (1999) no. 3, p. 91-93 | MR | Zbl | DOI
\bibitem{SS}J. Schmeling \& R. Siegmund-Schultze - “Hölder continuity of the holonomy maps for hyperbolic basic sets. I”, in Ergodic theory and related topics, III (Güstrow, 1990), Lect. Notes in Math., vol. 1514, Springer, Berlin, 1992, p. 174-191 | Zbl | DOI
\bibitem{Sigmund}K. Sigmund - “On minimal centers of attraction and generic points”, J. reine angew. Math. 295 (1977), p. 72-79 | MR | Zbl | DOI
\bibitem{Smale}S. Smale - “Differentiable dynamical systems”, Bull. Amer. Math. Soc. 73 (1967), p. 747-817 | MR | Zbl | DOI
\bibitem{Stewart}G. W. Stewart - Matrix algorithms. Vol. I. Basic decompositions, Society for Industrial and Applied Mathematics, Philadelphia, PA, 1998 | Zbl | DOI
\bibitem{Viana}M. Viana - “Almost all cocycles over any hyperbolic system have nonvanishing Lyapunov exponents”, Ann. of Math. (2) 167 (2008) no. 2, p. 643-680 | MR | Zbl | DOI
\bibitem{Walters}P. Walters - An introduction to ergodic theory, Graduate Texts in Math., vol. 79, Springer-Verlag, New York-Berlin, 1982 | MR | Zbl
\bibitem{Wirth}F. Wirth - “The generalized spectral radius and extremal norms”, Linear Algebra Appl. 342 (2002), p. 17-40 | MR | Zbl | DOI
\end{thebibliography}
